\subsection{Respuesta técnica y plazos de seguridad}\label{sec:sd8-respuesta-seguridad}
Estas reglas complementan AC-R-09/10/13 y las acciones de IN3, conservando ejecutor, propietario y paquete de cada ficha. Seguridad dirige evaluación/contención; SRE ejecuta aislamiento y restitución autorizados; Datos preserva evidencia y concilia; Calidad verifica corrección y negocio. Proyecto coordina comunicación y capacidad. El catálogo de escenarios de 8.2 remite a SEC-1.0/IR-1.0, SEC-ROT-1.0/GLS-1.0, DEP-6, AUD-5 e IN3-C01--C05; no se declara eficacia por escribir estas instrucciones.

El programa de vulnerabilidades mantiene activo, hallazgo, severidad, fecha de publicación/detección, responsable, vencimiento, corrección y prueba. El máximo de remediación es siete días corridos para críticas, quince para altas y treinta para medias, desde publicación o detección, sin reiniciar reloj al asignar técnico ni confundirlo con SLA de incidente. Si no existe parche, se retira, sustituye o deshabilita la función afectada y se comprueba remediación; control compensatorio sin evidencia o congelamiento no extienden el plazo. Se aplica regresión y retorno antes de reabrir, conservando hechos y evidencia. \parencites[art.~21.1]{sd8-ba}[RT-11.04]{sd8-bt}

IR-1.0 abre cronología desde detección y escala inmediatamente el crítico. El CLIENTE recibe comunicación de incidente crítico dentro de dos horas desde detección; toda brecha de seguridad o datos personales se notifica dentro de veinticuatro horas con informe preliminar y análisis de causa raíz dentro de los cinco días hábiles siguientes. Un incidente crítico que además sea brecha cumple ambos avisos: no espera al de veinticuatro horas para el primero. Seguridad/SRE preservan original, alcance conocido, incertidumbre, medidas, responsables y próxima actualización; Proyecto coordina los canales funcionales con la Contraparte, sin esperar comité ni cierre de investigación. \parencites[art.~21.3]{sd8-ba}[RT-11.18/19]{sd8-bt}

Ransomware o fuga requieren aislar el activo/frontera afectado, revocar credencial, custodiar evidencia y comprobar copias/claves antes de restaurar; la captura local segura se preserva sin aislar simultáneamente el par por defecto. DDoS o saturación limitan tráfico diferible y aplican capacidad/prioridad de SD4 sin descartar registros críticos. Dispositivo comprometido pierde facultades y sincroniza hechos sólo tras comprobación; una orden remota durante WAN caída no se presenta como ejecutada. Pérdida de observabilidad mantiene incidente abierto y escalamiento alternativo hasta comprobar detección y respuesta, sin declarar cobertura por contar un puesto SOC nominal.

La habilitación o rehabilitación de IN3 exige comprobar IN3-C01--C05 y el modelado propio: no escritura sin revisión, fragmento/versión trazables, permiso y tratamiento autorizado, resultados por campo/formato, desactivación e ingreso manual con conciliación de respuestas tardías. Región, versión/configuración y garantía de no entrenamiento permanecen puerta previa a tratar documentos del CLIENTE. Proyecto no consume una reserva distinta por cada escenario del mismo incidente; recursos, pruebas y suplencias siguen sujetos a la conciliación de T-15.
